% Zone „Das kennst du schon“ – aus bank/einheiten/zone.jsonl (zusammenbau v0.1)
\einheitenkopf[zone][Kennst du schon]{Das kennst du schon}
\setcounter{aufgabe}{0}
%% TODO Grundvorstellungs-Aufgabe als erste Hauptnummer der Zone (2.8) fehlt in der Bank

%% TODO Titel: Ich-kann-Satz fehlt in der Bank (Platzhalter: Kettenname) – Nr. 1
%% TODO Anweisung: ein Satz über der ersten Teilaufgabe fehlt in der Bank – Nr. 1
\begin{aufgabe}{Multiplizieren und Dividieren mit zehn, hundert und tausend als Kommaverschiebung, auch bei Dezimalzahlen}
%% TODO Darstellung für schwach fehlt in der Bank (einheiten-zone-f1-v1; Abschnitt „Für schwache Schüler“)
\swz{$4{,}2$ mal hundert?}{}
%% TODO Darstellung für schwach fehlt in der Bank (einheiten-zone-f1-v3; Abschnitt „Für schwache Schüler“)
\swz{$0{,}36$ mal tausend?}{}
\end{aufgabe}

%% TODO Titel: Ich-kann-Satz fehlt in der Bank (Platzhalter: Kettenname) – Nr. 2
%% TODO Anweisung: ein Satz über der ersten Teilaufgabe fehlt in der Bank – Nr. 2
\begin{aufgabe}{Stellenwerttafel lesen und schreiben, Nullen an der richtigen Stelle ergänzen (null Komma null sechs Meter; eintausendfünfhundert Milliliter)}
%% TODO Darstellung für schwach fehlt in der Bank (einheiten-zone-f2-v1; Abschnitt „Für schwache Schüler“)
\swz{$3$ Einer, $4$ Zehntel, $7$ Hundertstel – welche Zahl?}{}
%% TODO Darstellung für schwach fehlt in der Bank (einheiten-zone-f2-v3; Abschnitt „Für schwache Schüler“)
\swz{$16$ Hundertstel als Dezimalzahl?}{}
\end{aufgabe}

%% TODO Titel: Ich-kann-Satz fehlt in der Bank (Platzhalter: Kettenname) – Nr. 3
%% TODO Anweisung: ein Satz über der ersten Teilaufgabe fehlt in der Bank – Nr. 3
\begin{aufgabe}{Dezimalzahlen mit einer Zahl malnehmen und durch eine Zahl teilen, auch schriftlich oder mit dem Taschenrechner}
%% TODO Darstellung für schwach fehlt in der Bank (einheiten-zone-f3-v1; Abschnitt „Für schwache Schüler“)
\swz{$0{,}3 \cdot 4$?}{}
%% TODO Darstellung für schwach fehlt in der Bank (einheiten-zone-f3-v3; Abschnitt „Für schwache Schüler“)
\swz{$1{,}25 \cdot 6$?}{}
\end{aufgabe}

%% TODO Titel: Ich-kann-Satz fehlt in der Bank (Platzhalter: Kettenname) – Nr. 4
%% TODO Anweisung: ein Satz über der ersten Teilaufgabe fehlt in der Bank – Nr. 4
\begin{aufgabe}{Dezimalzahlen vergleichen und ordnen (stellenweise, Nullen anhängen)}
%% TODO Darstellung für schwach fehlt in der Bank (einheiten-zone-f4-v1; Abschnitt „Für schwache Schüler“)
\swz{Größere Zahl: $0{,}7$ oder $0{,}9$?}{}
%% TODO Darstellung für schwach fehlt in der Bank (einheiten-zone-f4-v3; Abschnitt „Für schwache Schüler“)
\swz{Ordne, mit der kleinsten beginnend: $0{,}45$; $0{,}405$; $0{,}54$}{}
\end{aufgabe}

%% TODO Titel: Ich-kann-Satz fehlt in der Bank (Platzhalter: Kettenname) – Nr. 5
%% TODO Anweisung: ein Satz über der ersten Teilaufgabe fehlt in der Bank – Nr. 5
\begin{aufgabe}{Skala am Lineal und an der Uhr lesen (Untereinheiten)}
\swz{Welche Zahl gehört zu $A$, welche zu $B$?}{\zahlenstrahl[xmin=0,xmax=3,xstep=0.1,karo=0.5]{0.7/A, 2.3/B}}
\swz{Welche Zahl gehört zu $C$?}{\zahlenstrahl[xmin=4,xmax=6,xstep=0.1,karo=0.6]{5.4/C}}
\end{aufgabe}

%% TODO Titel: Ich-kann-Satz fehlt in der Bank (Platzhalter: Kettenname) – Nr. 6
%% TODO Anweisung: ein Satz über der ersten Teilaufgabe fehlt in der Bank – Nr. 6
\begin{aufgabe}{Längen im Koordinatensystem als Differenz von Koordinaten lesen (Abstand zweier Schnittstellen, Höhe zwischen Tiefpunkt und Scheitel, Funktionswert als Höhe), Flächeninhalte in Flächeneinheiten aus dem Integral oder der Figur und Volumina aus Querschnitt mal Länge}
%% TODO Darstellung für schwach fehlt in der Bank (einheiten-zone-f6-v1; Abschnitt „Für schwache Schüler“)
\swz{Punkte $P(2|1)$ und $Q(2|6)$ – senkrechter Abstand?}{}
%% TODO Darstellung für schwach fehlt in der Bank (einheiten-zone-f6-v3; Abschnitt „Für schwache Schüler“)
\swz{Tiefpunkt $T(1|-1{,}5)$, Scheitel $S(0|2{,}25)$ – senkrechter Abstand?}{}
\end{aufgabe}

%% TODO Titel: Ich-kann-Satz fehlt in der Bank (Platzhalter: Kettenname) – Nr. 7
%% TODO Anweisung: ein Satz über der ersten Teilaufgabe fehlt in der Bank – Nr. 7
\begin{aufgabe}{Bruchteil einer Größe bilden (ein Viertel von einem Kilogramm; die Hälfte einer Stunde)}
%% TODO Darstellung für schwach fehlt in der Bank (einheiten-zone-f7-v1; Abschnitt „Für schwache Schüler“)
\swz{Ein Viertel von $8$ kg?}{}
%% TODO Darstellung für schwach fehlt in der Bank (einheiten-zone-f7-v3; Abschnitt „Für schwache Schüler“)
\swz{Drei Viertel von $12$ l?}{}
\end{aufgabe}

%% TODO Titel: Ich-kann-Satz fehlt in der Bank (Platzhalter: Kettenname) – Nr. 8
%% TODO Anweisung: ein Satz über der ersten Teilaufgabe fehlt in der Bank – Nr. 8
\begin{aufgabe}{Quadrat und dritte Potenz einer Länge mit Einheit (zwanzig Zentimeter zum Quadrat sind vierhundert Quadratzentimeter)}
%% TODO Darstellung für schwach fehlt in der Bank (einheiten-zone-f8-v1; Abschnitt „Für schwache Schüler“)
\swz{$3$ cm zum Quadrat – wie viel cm²?}{}
%% TODO Darstellung für schwach fehlt in der Bank (einheiten-zone-f8-v3; Abschnitt „Für schwache Schüler“)
\swz{$0{,}4$ m hoch drei – wie viel m³?}{}
\end{aufgabe}

%% TODO Titel: Ich-kann-Satz fehlt in der Bank (Platzhalter: Kettenname) – Nr. 9
%% TODO Anweisung: ein Satz über der ersten Teilaufgabe fehlt in der Bank – Nr. 9
\begin{aufgabe}{Prozentwert einer Größe berechnen (zehn Prozent eines Volumens)}
%% TODO Darstellung für schwach fehlt in der Bank (einheiten-zone-f9-v1; Abschnitt „Für schwache Schüler“)
\swz{$25\,\%$ von $36$ kg?}{}
%% TODO Darstellung für schwach fehlt in der Bank (einheiten-zone-f9-v3; Abschnitt „Für schwache Schüler“)
\swz{$15\,\%$ von $2{,}4$ m³?}{}
\end{aufgabe}

%% TODO Titel: Ich-kann-Satz fehlt in der Bank (Platzhalter: Kettenname) – Nr. 10
%% TODO Anweisung: ein Satz über der ersten Teilaufgabe fehlt in der Bank – Nr. 10
\begin{aufgabe}{Formel nach einer Größe umstellen (aus einer Verhältnisgleichung die gesuchte Größe isolieren)}
%% TODO Darstellung für schwach fehlt in der Bank (einheiten-zone-f10-v1; Abschnitt „Für schwache Schüler“)
\swz{Löse: $x \cdot 4 = 28$}{}
%% TODO Darstellung für schwach fehlt in der Bank (einheiten-zone-f10-v3; Abschnitt „Für schwache Schüler“)
\swz{Löse: $1{,}6 \cdot x = 12$}{}
\end{aufgabe}

%% TODO Titel: Ich-kann-Satz fehlt in der Bank (Platzhalter: Kettenname) – Nr. 11
%% TODO Anweisung: ein Satz über der ersten Teilaufgabe fehlt in der Bank – Nr. 11
\begin{aufgabe}{Fläche eines Rechtecks und eines Dreiecks berechnen (für den Materialbedarf)}
%% TODO Darstellung für schwach fehlt in der Bank (einheiten-zone-f11-v1; Abschnitt „Für schwache Schüler“)
\swz{Rechteck, $6$ m lang und $4$ m breit – Fläche?}{}
%% TODO Darstellung für schwach fehlt in der Bank (einheiten-zone-f11-v3; Abschnitt „Für schwache Schüler“)
\swz{Rechteck, $2{,}4$ m lang und $1{,}5$ m breit – Fläche?}{}
\end{aufgabe}

%% TODO Titel: Ich-kann-Satz fehlt in der Bank (Platzhalter: Kettenname) – Nr. 12
%% TODO Anweisung: ein Satz über der ersten Teilaufgabe fehlt in der Bank – Nr. 12
\begin{aufgabe}{Multiplizieren und Dividieren mit zehn, hundert und tausend als Kommaverschiebung, auch bei Dezimalzahlen – Fehler finden}
\swz[3]{Lena rechnet: $8{,}4$ geteilt durch hundert ergibt $840$. Wo ist der Fehler?}{}
\end{aufgabe}

%% TODO Titel: Ich-kann-Satz fehlt in der Bank (Platzhalter: Kettenname) – Nr. 13
%% TODO Anweisung: ein Satz über der ersten Teilaufgabe fehlt in der Bank – Nr. 13
\begin{aufgabe}{Multiplizieren und Dividieren mit zehn, hundert und tausend als Kommaverschiebung, auch bei Dezimalzahlen – selbst rechnen}
%% TODO Darstellung für schwach fehlt in der Bank (einheiten-zone-f1-v6; Abschnitt „Für schwache Schüler“)
\swz{$6{,}5$ geteilt durch hundert?}{}
\end{aufgabe}
